% !PRINTED 323
\weberpage{323}

If $D$ is a discriminant and $n$ a rational integer, then we define a symbol $(D,n)$ by the following stipulations:
\weq{1}{(D,0)=0.}
\weq{2}{(D,1)=1.}
\weq{3}{\begin{aligned} (D,-1) &= +1, & D>0.\\ &= -1, & D<0. \end{aligned}}
\weq{4}{(D,p)=0,}
if $p$ is a prime divisor of $D$.
\weq{5}{(D,2)=(-1)^{\frac{D^2-1}{8}},\quad D\ \text{odd}.}
\weq{6}{(D,p)=\left(\dfrac{D}{p}\right),}
if $p$ is an odd prime that does not divide $D$.

If one decomposes $\pm n$ into its prime factors
\[
n = \pm p,\ p',\ p'',\ \dots,
\]
then
\weq{7}{(D,n) = (D,\pm 1)(D,p)(D,p')(D,p'')\dots}

By (1) to (7) the symbol $(D,n)$ is evidently defined without contradiction for all numbers $n$, and can have only one of the values $-1$, $0$, $+1$. It is always $=0$ then and only then, when $D$ and $n$ are not relatively prime to one another. From (7) there follows further:
\weq{8}{(D,n)(D,n') = (D,nn'),}
and if $Q^2$ is a square number that has no common divisor with $n$,
\weq{9}{(Q^2D,n) = (D,n).}

The formulae relating to the law of reciprocity take a particular form in this notation. We consider first the \gesperrt{prime discriminants}.

If $\pm p \equiv 1$ (mod 4), then by (3), (5), (6):
\[
(\pm p, -1) = (-1)^{\frac{p-1}{4}} = \left(\dfrac{-1}{p}\right),
\]
\[
(\pm p, 2) = (-1)^{\frac{p^2-1}{8}} = \left(\dfrac{2}{p}\right),
\]
\[
(\pm p, n) = \left(\dfrac{\pm p}{n}\right) = \left(\dfrac{n}{p}\right),
\]
%CONT>>
